\section{Convergence theorems for expectations}

We want to study notions of continuity for the
expectation operator. If $X_n\to X$ as $n\to\infty$,
under what conditions do we have that
$\mathbf{E}(X_n)\to \mathbf{E} X$? First we have to
decide what ``$X_n\to X$'' actually \emph{means}. We
define two notions of convergence of a sequence of
random variables to a limit.

\begin{defn}
Let $X, X_1, X_2, \ldots$ be random variables all
defined on the same probability space. We say that
\textbf{$X_n$ converges in probability to $X$}, and denote
$X_n \xrightarrow[n\to\infty]{\mathrm{P}} X$, if for
all $\epsilon>0$ we have that

\[
\mathrm{P}(|X_n-X|>\epsilon) \xrightarrow[n\to\infty]{} 0.
\]

\end{defn}

\begin{defn}
With $X,X_1,X_2,\ldots$ as before, we say that \textbf{$X_n$
converges almost surely to $X$} (or \textbf{converges to $X$
with probability 1}), and denote
$X_n \xrightarrow[n\to\infty]{\textrm{a.s.}} X$, if

\[
\mathrm{P}( X_n \to X ) = \mathrm{P}\left(\left\{ \omega \in \Omega : X(\omega)=\lim_{n\to\infty} X_n(\omega) \right\} \right) = 1.
\]

\end{defn}

\begin{xexercise}
Show that
$\left\{ \omega \in \Omega :
X(\omega)=\lim_{n\to\infty} X_n(\omega) \right\}$
is an event and therefore has a well-defined
probability. In other words, represent it in terms of
countable union, intersection and complementation
operations on simple sets that are known to be events.
Hint: Use the $\epsilon-\delta$ definition of a
limit.
\end{xexercise}

\begin{lem}
\label{lem-convergence-types}
Almost sure convergence is a stronger notion of
convergence than convergence in probability. In other
words, if
$X_n\xrightarrow[n\to\infty]{\textrm{a.s.}} X$ then
$X_n\xrightarrow[n\to\infty]{\mathrm{P}} X$, but the
converse is not true.
\end{lem}

\begin{xexercise}
Prove \hyperref[lem-convergence-types]{Lemma~\ref*{lem-convergence-types}}. For the
counterexample showing that convergence in probability
does not imply almost sure convergence, consider the
following sequence of random variables defined on the
space $((0,1),{\cal B}, \textrm{Lebesgue measure})$:

\begin{center}
\fitwidth{$\displaystyle
\begin{array}{l}
1_{(0,1)}, \\
1_{(0,1/2)},  1_{(1/2,1)}, \\
1_{(0,1/4)},  1_{(1/4,2/4)},  1_{(2/4,3/4)},  1_{(3/4,1)},  \\
1_{(0,1/8)},  1_{(1/8,2/8)},  1_{(2/8,3/8)},  1_{(3/8,4/8)}, 1_{(4/8,5/8)},  1_{(5/8,6/8)},  1_{(6/8,7/8)},  1_{(7/8,1)}, \\ \ldots
\end{array}
$}
\end{center}

\end{xexercise}

\begin{lem}
\label{lem:prob-as-subseq}
If $(X_n)_{n=1}^\infty$ is a sequence of r.v.s such
that $X_n \xrightarrow[n\to\infty]{\mathrm{P}} X$
then there exists a subsequence
$(X_{n_k})_{k=1}^\infty$ such that
$X_n \xrightarrow[k\to\infty]{\textrm{a.s.}} X$.
\end{lem}

\begin{xexercise}
Prove \hyperref[lem:prob-as-subseq]{Lemma~\ref*{lem:prob-as-subseq}}.
\end{xexercise}

\noindent We can now formulate the fundamental convergence
theorems for Lebesgue integration.

\begin{thm}[Bounded convergence theorem]
If $X_n$ is a sequence of r.v.'s such that
$|X_n|\le M$ for all $n$, and $X_n\to X$ in
probability, then $\mathbf{E} X_n \to \mathbf{E} X$.
\end{thm}

\begin{proof}
Fix $\epsilon>0$. Then

\begin{eqnarray*}
|\mathbf{E} X_n - \mathbf{E} X| &\le& \mathbf{E} |X_n - X| = \mathbf{E} |X_n - X|1_{\{|X_n-X|>\epsilon\}} + \mathbf{E} |X_n - X|1_{\{|X_n-X|\le \epsilon\}} \\
&\le& 2M \mathrm{P}(|X_n-X|>\epsilon) + \epsilon\xrightarrow[n\to\infty]{} \epsilon.
\end{eqnarray*}

\noindent Since $\epsilon$ was an arbitrary positive number,
this implies that
$|\mathbf{E} X_n-\mathbf{E} X|\to 0$, as claimed.
\end{proof}

\begin{thm}[Fatou's lemma]
If $X_n \ge 0$ then
$\liminf_{n\to\infty} \mathbf{E} X_n \ge
\mathbf{E}(\liminf_{n\to\infty} X_n).$
\end{thm}

\noindent To see that the inequality in the lemma can fail to
be an equality, let $U \sim U(0,1)$, and define
$X_n = n 1_{\{U\le 1/n\}}$. Clearly
$\liminf_{n\to\infty}X_n=\lim_{n\to\infty} X_n \equiv
0$,
but $\mathbf{E}(X_n) = 1$ for all $n$.

\begin{proof}
Let $Y = \liminf_{n\to\infty} X_n$. Note that $Y$ can
be written as

\[
Y = \sup_{n\ge 1} \inf_{m\ge n} X_m
\]

\noindent (this is a general fact about the lim inf of a
sequence of real numbers), or $Y = \sup_n Y_n$, where
we denote

\[
Y_n = \inf_{m\ge n} X_m.
\]

\noindent We have $Y_n \le X_n$, and as $n\to\infty$,
$Y_n \to Y$ (in fact $Y_n \uparrow Y$) almost surely.
Therefore $\mathbf{E} Y_n \le \mathbf{E} X_n$, so
$\liminf_{n\to\infty} \mathbf{E} Y_n \le
\liminf_{n\to\infty} \mathbf{E} X_n$,
and therefore it is enough to show that

\[
\liminf_{n\to\infty} \mathbf{E} Y_n \ge \mathbf{E} Y.
\]

\noindent But for any $M$ we have that

\[
Y_n \wedge M \xrightarrow[n\to\infty]{\textrm{a.s.}} Y \wedge M,
\]

\noindent and this is a sequence of uniformly bounded r.v.'s,
therefore by the bounded convergence theorem we get
that

\[
\mathbf{E}(Y_n) \ge \mathbf{E}(Y_n \wedge M) \xrightarrow[n\to\infty]{} \mathbf{E}(Y\wedge M).
\]

\noindent We therefore get that
$\liminf_{n\to\infty} \mathbf{E}(Y_n)\ge
\mathbf{E}(Y\wedge M)$
for any $M>0$, which implies the result because of
the following exercise.
\end{proof}

\begin{xexercise}
Let $Y\ge 0$ be a random variable. Prove that

\[
\mathbf{E}(Y) = \sup_{M>0} \mathbf{E}(Y\wedge M).
\]

\end{xexercise}

\begin{thm}[Monotone convergence theorem]
If $0\le X_n \uparrow X$ as $n\to\infty$ then
$\mathbf{E} X_n \uparrow \mathbf{E} X$.
\end{thm}

\begin{proof}

\[
\mathbf{E} X = \mathbf{E} [\liminf_{n\to\infty} X_n] \le \liminf_{n\to\infty} \mathbf{E} X_n \le \limsup_{n\to\infty} \mathbf{E} X_n \le \limsup_{n\to\infty}  \mathbf{E} X = \mathbf{E} X.
\]

\end{proof}

\begin{thm}[Dominated convergence theorem]
If $X_n\to X$ almost surely, $|X_n|\le Y$ for all
$n\ge 1$ and $\mathbf{E} Y<\infty$, then
$\mathbf{E} X_n\to \mathbf{E} X$.
\end{thm}

\begin{proof}
Apply Fatou's lemma separately to $Y+X_n$ and to
$Y-X_n$.
\end{proof}
